\documentclass[a4paper, 12pt, oneside]{book}
\usepackage[subpreambles=true]{standalone}
\usepackage{import}

\usepackage[a4paper, width=150mm, top=25mm, bottom=25mm, headheight=10mm]{geometry}
% \usepackage[a4paper, width=150mm, top=25mm, bottom=25mm, headheight=10mm, bindingoffset=6mm]{geometry}

\usepackage{lmodern}

\usepackage{setspace}
\onehalfspacing

\usepackage{fancyhdr}
\renewcommand{\headrulewidth}{0pt}
\renewcommand{\footrulewidth}{1.2pt}
\pagestyle{fancy}
\fancyhead{}
\fancyhead[R]{\leftmark}
\fancyfoot{}
\fancyfoot[R]{page \thepage}
\newcommand{\footnotecustom}{}
\fancypagestyle{plain}{
	\fancyhead{}
	\fancyfoot{}
	\fancyfoot[R]{page \thepage}
	\fancyfoot[L]{\footnotecustom}
}

\usepackage{mhchem}
\usepackage{siunitx}

\begin{document}
\begin{center}
	\Large \bfseries Electromagnetic Contribution to Nuclear Mass
	\hline
\end{center}

Consider two mirror nuclei, \ce{^15_7N} and \ce{^15_8O}. Since \ce{^15_7N} has more neutrons, a
naive expectation would be to assume \ce{^15_7N} has more mass as neutrons are more massive compared
to protons. However, this is not the case, \ce{^15_8O} is the more massive nucleus. In this note, we
will argue that electrostatic potential, through mass energy correspondence, can account for this
discrepancy to some degree.

A uniformly charged sphere of radius $R$ and total charge $Q$ has the following potential energy.
\begin{equation}
	V = \frac{3}{20 \pi \epsilon_0} \, \frac{Q^2}{R}
\end{equation}
Let, $R$ be the effective radius of a nucleus and $Z$ its number of protons. Our goal is to
calculate the binding energy due to electric potential. We do it in the following way. First $Z$
protons are considered in their ``free" state, i.e., $Z$ charged spheres of radius $R$ that sit
infinitely far apart from each other such that they do not interact with each other. Then, the
spheres are brought together and a new charged sphere is created with charge $Ze$ and radius $R$. The
difference in the potential energy between these two configurations, $U$, is assumed to contribute
to mass through $E=mc^2$.
\begin{equation}
	U = \frac{3}{20 \pi \epsilon_0} \, \frac{(Ze)^2 - Z e^2}{R}
\end{equation}
If $N$ denotes the number of neutrons in a nuclei, then a nuclear configuration is denoted by the
number of protons and neutrons as $(Z, N)$. Now, $(Z, N)$ and $(N, Z)$ are mirror nuclei of each
other. We assume that the binding energy due to the nuclear forces are equal in both nuclei as
nuclear force treats both proton and neutron almost equally. The mass difference between mirror
nuclei then comes from the difference in electric potential energy and due to the mass difference
between proton and neutron. The mass difference between $(N, Z)$ and $(Z, N)$ then simplifies to the
following expression.
\begin{equation}
	\Delta m = \frac{3e^2}{20 \pi \epsilon_0 c^2} \frac{(N-Z)(N+Z-1)}{R_\mathrm{avg}} - (N-Z) (m_n - m_p)
\end{equation}
We assumed both mirror nuclei have almost the same radius and $R_\mathrm{avg}$ is their average.
$m_n, m_p$ are neutron and proton masses respectively. If expressed in electron volts the mass
difference becomes the following.
\begin{equation}
	\Delta m = \frac{3e}{20 \pi \epsilon_0} \frac{(N-Z)(N+Z-1)}{R_\mathrm{avg}} - (N-Z) (m_n - m_p)
\end{equation}
For \ce{^15_7N}, $Z=7, N=8$. The r.m.s. radius is $2.61\unit{\femto\meter}$, which we have to
multiply by $\sqrt{5/3}$ to get $R_\mathrm{avg}$ and $m_n - m_p$ is
$1.293\unit{\mega\electronvolt}$. The mass difference then comes out to be
$2.29\unit{\mega\electronvolt}$, as in \ce{^17O} has $2.29\unit{\mega\electronvolt}$ more mass than
\ce{^15N}. The experimental value for atomic mass difference in $2.754\unit{\mega\electronvolt}$,
but Oxygen has an extra electron, so we need to subtract the energy of an electron. Then, the
nuclear mass difference is $2.24\unit{\mega\electronvolt}$, very close to the value given by the
formula. If we only used the mass difference between proton and neutron we would guess that
\ce{^15N} is $1.293\unit{\mega\electronvolt}$ more massive than \ce{^15O}. But, due to Oxygen having
more protons, the increase in electric potential (scales as $\sim Z^2$) more than compensates for
the lower mass of proton.

\end{document}
